% !TeX program = pdflatex
% Pakiet dokumentów do aktu zawiązania Basilisk Systems P.S.A. — uchwały założycielskie, pierwsze uchwały Rady Dyrektorów, regulamin Walnego Zgromadzenia.
% Odesłania do umowy spółki pobierane z psa.aux (xr-hyper).
\documentclass[11pt,a4paper]{report}
\usepackage[T1]{fontenc}
\usepackage[utf8]{inputenc}
\usepackage[english,polish]{babel}
\usepackage{xr-hyper}
\usepackage{basilisk_i18n}
\usepackage{lmodern}
\usepackage{microtype}
\usepackage{setspace}
\usepackage{parskip}
\usepackage{needspace}
\externaldocument{psa}[psa.pdf]

\BusinessPlanCompany{Basilisk Systems}
\BusinessPlanWatermarkText{Basilisk Systems}
\BusinessPlanWatermarkColor{black!6}
\BusinessPlanWatermarkFontSize{38}{46}
\BusinessPlanWatermarkAngle{38}
\BusinessPlanAuthorsText{Basilisk Systems --- Założyciele}
\BusinessPlanConfidentialText{Projekty uchwał --- poufne}
\BusinessPlanFooterText{Basilisk Systems --- pakiet dokumentów do aktu zawiązania (projekt)}
\BusinessPlanVersionText{Wersja 1.0 --- projekt}

\setstretch{1.12}
\setlist[enumerate]{topsep=0.2em,itemsep=0.15em,parsep=0.05em,leftmargin=2.0em}
\setlength{\parindent}{0pt}
\setlength{\parskip}{0.5em}
\setlength{\emergencystretch}{4em}

\newcounter{paragraf}
\newcommand{\paragraf}[1]{%
  \refstepcounter{paragraf}%
  \Needspace{5\baselineskip}%
  \subsection*{\S~\theparagraf\ #1}%
}
\newenvironment{ContractClause}[1]{\paragraf{#1}}{\par}
\newcommand{\field}[1]{\textbf{[#1]}}
\newcommand{\CompanyFull}{Basilisk Systems prosta spółka akcyjna}
\newcommand{\US}[1]{§~#1 umowy spółki}
\newcommand{\Uchwala}[2]{\Needspace{8\baselineskip}\section*{#1}\addcontentsline{toc}{section}{#1 --- #2}\textbf{w sprawie: #2}\par}

\hypersetup{
  pdftitle={Basilisk Systems P.S.A. --- pakiet dokumentów do aktu zawiązania},
  pdfauthor={Basilisk Systems}
}

\begin{document}

\BusinessPlanTitlePage
  {Basilisk Systems}
  {Wersja 1.0 --- projekt}
  {Pakiet aktu zawiązania}
  {Uchwały założycielskie, pierwsze uchwały Rady Dyrektorów i regulamin Walnego Zgromadzenia}
  {Do umowy spółki Basilisk Systems P.S.A. w wersji 3.3-S}
  {1.0}

\tableofcontents
\newpage

\chapter*{Wprowadzenie}
\addcontentsline{toc}{chapter}{Wprowadzenie}

Umowa spółki wymaga, aby Założyciele w uchwałach objętych aktem zawiązania Spółki wskazali pierwszych dyrektorów, podmiot prowadzący rejestr akcjonariuszy, pierwszy regulamin Walnego Zgromadzenia oraz pełnomocnika do czynności z Założycielami będącymi dyrektorami (\US{\ref{art:final} ust.~6}). Niniejszy pakiet zawiera projekty tych uchwał, a także pierwszych uchwał Rady Dyrektorów podejmowanych bezpośrednio po zawiązaniu.

\textbf{Kolejność czynności w dniu zawiązania:}
\begin{enumerate}
  \item podpisanie umowy spółki w formie aktu notarialnego wraz z uchwałami nr 1--5 (część~I);
  \item niezwłocznie po podpisaniu --- zawarcie przez pełnomocnika z uchwały nr 4 umów przeniesienia praw stanowiących wkłady niepieniężne (Umowa nr~1) oraz podpisanie umowy akcjonariuszy;
  \item posiedzenie Rady Dyrektorów i uchwały nr R1--R3 (część~II);
  \item wniesienie wkładów pieniężnych i zgłoszenie Spółki do rejestru przedsiębiorców.
\end{enumerate}

Bez uchwały nr 4 umowy zawierane przez Spółkę z Założycielami będącymi dyrektorami są zagrożone nieważnością (art.~300$^{79}$ §~1 KSH). Pełnomocnikiem nie powinna być strona danej umowy ani inny dyrektor.

% ================================================================
\chapter*{Część I --- Uchwały Założycieli objęte aktem zawiązania}
\addcontentsline{toc}{chapter}{Część I --- Uchwały Założycieli objęte aktem zawiązania}

\Uchwala{Uchwała nr 1}{powołania pierwszych dyrektorów}
Działając na podstawie \US{\ref{art:rada} ust.~2--3} oraz \US{\ref{art:final} ust.~6}, Założyciele powołują do Rady Dyrektorów na pierwszą wspólną kadencję:
\begin{enumerate}
  \item \field{imię i nazwisko} --- dyrektor;
  \item \field{imię i nazwisko} --- dyrektor;
  \item \field{imię i nazwisko} --- dyrektor niewykonawczy.
\end{enumerate}
Rada Dyrektorów liczy \field{3--7} dyrektorów. Co najmniej jeden dyrektor ma być dyrektorem niewykonawczym (\US{\ref{art:rada} ust.~4}). Uchwała wchodzi w życie z chwilą podjęcia.

\Uchwala{Uchwała nr 2}{wyboru podmiotu prowadzącego rejestr akcjonariuszy}
Działając na podstawie art.~300$^{31}$ KSH oraz \US{\ref{art:akcje} ust.~5}, Założyciele wybierają \field{firma podmiotu uprawnionego, siedziba, KRS} jako podmiot prowadzący rejestr akcjonariuszy Spółki i upoważniają Radę Dyrektorów do zawarcia z nim umowy. Rada Dyrektorów zgłasza ten podmiot do rejestru przedsiębiorców w terminach wynikających z przepisów prawa.

\Uchwala{Uchwała nr 3}{przyjęcia regulaminu Walnego Zgromadzenia}
Działając na podstawie \US{\ref{art:wz} ust.~5}, Założyciele przyjmują regulamin Walnego Zgromadzenia w brzmieniu stanowiącym część~III niniejszego pakietu.

\Uchwala{Uchwała nr 4}{powołania pełnomocnika do czynności z dyrektorami}
\begin{enumerate}
  \item Działając na podstawie art.~300$^{79}$ §~1 KSH oraz \US{\ref{art:reprezentacja} ust.~3} i \US{\ref{art:final} ust.~6}, Założyciele powołują \field{imię i nazwisko, PESEL albo numer dokumentu} na pełnomocnika Spółki do zawierania, zmiany, wykonywania i rozwiązywania w imieniu Spółki, z Założycielami będącymi dyrektorami, następujących umów oraz do dokonywania innych czynności z nimi związanych:
  \begin{enumerate}[label=\alph*)]
    \item umów przeniesienia praw stanowiących wkłady niepieniężne (Umowa nr~1, \US{\ref{art:cap-table} ust.~6}) wraz z protokołami wydania;
    \item umów przeniesienia praw do utworów tworzonych dla Spółki (Umowa nr~2, \US{\ref{art:ip} ust.~1}) oraz potwierdzeń kwartalnych;
    \item umowy akcjonariuszy wraz z jej załącznikami, w tym ofert, pełnomocnictw i oświadczeń o przystąpieniu;
    \item umów wykonawczych do obowiązku zbycia akcji (\US{\ref{art:obowiazek-zbycia} ust.~6}).
  \end{enumerate}
  \item Pełnomocnik nie może być stroną umowy, którą zawiera w imieniu Spółki. Pełnomocnictwo jest ważne do czasu jego odwołania uchwałą akcjonariuszy.
\end{enumerate}

\Uchwala{Uchwała nr 5}{zgody na czynności Spółki w organizacji}
Działając na podstawie \US{\ref{art:final} ust.~6}, Założyciele wyrażają zgodę na zaciągnięcie przez Spółkę w organizacji zobowiązań niezbędnych do jej zawiązania i rejestracji, w tym kosztów notarialnych, sądowych, rejestru akcjonariuszy i otwarcia rachunku bankowego, łącznie do kwoty \field{kwota} zł. Inne zobowiązania do dnia wpisu Spółki do rejestru przedsiębiorców wymagają zgody wszystkich Założycieli w formie dokumentowej.

\vspace{1em}
\field{podpisy wszystkich Założycieli --- w akcie notarialnym}

% ================================================================
\chapter*{Część II --- Pierwsze uchwały Rady Dyrektorów}
\addcontentsline{toc}{chapter}{Część II --- Pierwsze uchwały Rady Dyrektorów}

\Uchwala{Uchwała nr R1}{wyboru Przewodniczącego i wyznaczenia dyrektora wykonawczego}
Działając na podstawie \US{\ref{art:rada} ust.~4} i art.~300$^{76}$ KSH, Rada Dyrektorów:
\begin{enumerate}
  \item wybiera \field{imię i nazwisko} na Przewodniczącego Rady Dyrektorów;
  \item wyznacza \field{imię i nazwisko} na dyrektora wykonawczego i deleguje mu prowadzenie bieżącej działalności Spółki w granicach budżetu, z wyłączeniem spraw wymagających uchwały Rady Dyrektorów (\US{\ref{art:rada} ust.~7});
  \item stwierdza, że dyrektorem niewykonawczym jest \field{imię i nazwisko}.
\end{enumerate}

\Uchwala{Uchwała nr R2}{wyznaczenia osób odpowiedzialnych}
Rada Dyrektorów wyznacza:
\begin{enumerate}
  \item \field{imię i nazwisko} --- Odpowiedzialnego w rozumieniu polityki bezpieczeństwa informacji;
  \item \field{imię i nazwisko} --- kierownika przedsiębiorcy w rozumieniu ustawy o ochronie informacji niejawnych, na potrzeby przyszłego postępowania o wydanie świadectwa bezpieczeństwa przemysłowego;
  \item \field{imię i nazwisko} oraz \field{imię i nazwisko} --- dyrektorów odpowiedzialnych za przygotowanie wniosku o koncesję na wytwarzanie i obrót wyrobami i technologią o przeznaczeniu wojskowym.
\end{enumerate}

\Uchwala{Uchwała nr R3}{przyjęcia dokumentów wewnętrznych i stwierdzenia wniesienia wkładów}
\begin{enumerate}
  \item Rada Dyrektorów przyjmuje politykę bezpieczeństwa informacji, Cennik wynagrodzeń za rozwiązania projektowe oraz wzór umowy o zachowaniu poufności (\US{\ref{art:rada} ust.~7}).
  \item Rada Dyrektorów stwierdza wniesienie wkładów pieniężnych w łącznej kwocie \field{kwota} zł oraz --- po zawarciu umów przeniesienia i protokołów wydania --- wkładów niepieniężnych wymienionych w Załączniku nr~1 do umowy spółki, i ujawnia wynikającą z tego wysokość kapitału akcyjnego (\US{\ref{art:cap-table} ust.~5}).
\end{enumerate}

% ================================================================
\chapter*{Część III --- Regulamin Walnego Zgromadzenia}
\addcontentsline{toc}{chapter}{Część III --- Regulamin Walnego Zgromadzenia}

\begin{ContractClause}{Zakres}
Regulamin określa tryb zwoływania i obradowania Walnego Zgromadzenia oraz zasady udziału w nim przy wykorzystaniu środków komunikacji elektronicznej (\US{\ref{art:wz} ust.~5}; art.~300$^{92}$ KSH). Do uchwał podejmowanych poza Walnym Zgromadzeniem stosuje się \US{\ref{art:wz} ust.~3}.
\end{ContractClause}

\begin{ContractClause}{Zwołanie}
\begin{enumerate}
  \item Walne Zgromadzenie zwołuje Rada Dyrektorów zawiadomieniem wysłanym pocztą elektroniczną na adresy akcjonariuszy wpisane do rejestru akcjonariuszy, z zachowaniem terminów ustawowych (\US{\ref{art:wz} ust.~2 i 4}).
  \item Zawiadomienie wskazuje datę, godzinę, miejsce, porządek obrad, projekty uchwał, sposób udziału zdalnego i dane kontaktowe do zgłaszania problemów technicznych.
\end{enumerate}
\end{ContractClause}

\begin{ContractClause}{Udział zdalny}
\begin{enumerate}
  \item Akcjonariusz albo jego pełnomocnik może uczestniczyć w Walnym Zgromadzeniu przy wykorzystaniu środków komunikacji elektronicznej, w szczególności przez dwustronną komunikację audio-wideo w czasie rzeczywistym, oraz wykonywać prawo głosu na odległość.
  \item Identyfikacja następuje przez: a)~zgłoszenie udziału zdalnego pocztą elektroniczną z adresu wpisanego do rejestru akcjonariuszy najpóźniej na 24 godziny przed rozpoczęciem obrad; b)~okazanie dokumentu tożsamości podczas transmisji albo uwierzytelnienie kwalifikowanym podpisem elektronicznym; c)~w przypadku pełnomocnika --- przesłanie pełnomocnictwa w formie wymaganej przez prawo.
  \item Głos oddaje się ustnie do protokołu, przez funkcję głosowania platformy zapewniającą utrwalenie głosu albo wiadomością poczty elektronicznej z adresu wpisanego do rejestru akcjonariuszy, wysłaną w trakcie głosowania.
  \item Rada Dyrektorów zapewnia, że platforma umożliwia identyfikację uczestników, komunikację w czasie rzeczywistym i utrwalenie przebiegu obrad, oraz udostępnia ją akcjonariuszom co najmniej 30 minut przed rozpoczęciem obrad.
  \item Wymogi nałożone w niniejszym paragrafie nie mogą wykraczać poza niezbędne do identyfikacji akcjonariuszy i zapewnienia bezpieczeństwa komunikacji (art.~300$^{92}$ KSH).
\end{enumerate}
\end{ContractClause}

\begin{ContractClause}{Przebieg obrad}
\begin{enumerate}
  \item Walne Zgromadzenie otwiera Przewodniczący Rady Dyrektorów albo inny dyrektor. Spośród uprawnionych wybiera się przewodniczącego Walnego Zgromadzenia.
  \item Przewodniczący sporządza listę obecności, obejmującą uczestników zdalnych ze wskazaniem sposobu identyfikacji, stwierdza prawidłowość zwołania i zdolność do podejmowania uchwał oraz czuwa nad sprawnym przebiegiem obrad.
  \item Walne Zgromadzenie może zarządzić przerwę w obradach, w tym w razie zakłóceń technicznych dotyczących większej liczby uczestników. Problemy techniczne leżące po stronie akcjonariusza nie wpływają na ważność obrad i uchwał w zakresie dopuszczalnym przez prawo (\US{\ref{art:wz} ust.~6}).
  \item Uchwały, które z mocy prawa wymagają protokołu sporządzonego przez notariusza, w szczególności dotyczące zmiany umowy spółki, podejmuje się w obecności notariusza; notariusz uczestniczy w miejscu obrad.
\end{enumerate}
\end{ContractClause}

\begin{ContractClause}{Protokół i dane osobowe}
\begin{enumerate}
  \item Z obrad sporządza się protokół zawierający listę obecności, treść uchwał, liczbę głosów oddanych za każdą uchwałą i zgłoszone sprzeciwy. Protokół wpisuje się do księgi protokołów (art.~300$^{100}$ KSH).
  \item Nagranie obrad może być sporządzone wyłącznie w celu dokumentacji przebiegu Walnego Zgromadzenia, po uprzedzeniu uczestników, i jest przechowywane przez \field{okres}. Administratorem danych osobowych uczestników jest Spółka (art.~6 ust.~1 lit.~c i f RODO).
\end{enumerate}
\end{ContractClause}

\begin{ContractClause}{Postanowienia końcowe}
Zmiana regulaminu wymaga uchwały akcjonariuszy podjętej bezwzględną większością głosów oddanych. Regulamin wchodzi w życie z dniem zawiązania Spółki.
\end{ContractClause}

\end{document}
